\documentclass[11pt]{article}
\usepackage[utf8]{inputenc}
\usepackage[T1]{fontenc}
\usepackage{geometry}
\usepackage{hyperref}
\usepackage{enumitem}
\geometry{letterpaper, margin=1in}
\title{\textbf{Working Draft Outline: Paper 3 \\ Extended or Massive Constituent Route}}
\author{Stephen Breighner}
\date{March 2026}
\begin{document}
\maketitle

\section*{Status}
This outline covers the most speculative route: replace the standard massless graviton baseline with massive or extended bipolar constituents, then use internal structure to justify finite separation and short-range support.

\section*{Core Scope}
\begin{itemize}[leftmargin=1.25em]
  \item State clearly that this route departs from the standard GR graviton picture.
  \item Introduce the core-shell or extended-constituent idea as the source of short-range repulsion.
  \item Revisit whether the long-range potential remains Newtonian or becomes modified.
  \item Connect the constituent picture to the transfer-field branch if both are used together.
\end{itemize}

\section*{Proposed Flow}
\begin{enumerate}[leftmargin=1.4em]
  \item Motivate why a finite separation scale may require internal structure.
  \item Define the constituent assumptions: mass, shell, polarity, and overlap penalty.
  \item Show how those assumptions map onto the effective support term in the toy potential.
  \item Work through what must change in the large-scale force law and wave behavior.
  \item End with the consistency risks: extra modes, modified propagation, and derivation burden.
\end{enumerate}

\section*{Required Clarifications}
\begin{itemize}[leftmargin=1.25em]
  \item Does mass alone help, or is internal structure the actual support mechanism?
  \item What sets the shell scale or nonzero minimum separation?
  \item Which observations would most strongly constrain this route first?
\end{itemize}

\section*{Near-Term Deliverable}
A separate speculative paper that explicitly treats the constituent picture as a branch model rather than as an assumption hidden inside the first paper.

\end{document}
